\section{State: recurrent form, the HMM reading, and its limits}
\label{sec:state}

The complaint addressed here is that decoder-only transformers are ``just gated
RNNs'', and that constraining such an RNN should yield HMM-like decoding. We
find the first half true but nearly empty as usually stated
(Section~\ref{sec:state-form}), the HMM identity true but \emph{not} a statement
about linear attention (Section~\ref{sec:state-hmm}), and we prove that the
missing bridge cannot be built (Theorem~\ref{thm:noembed}). We close with the
capacity bound that says what unbounded state buys
(Theorem~\ref{thm:recall}).

\subsection{The recurrent form is real, and almost vacuous}
\label{sec:state-form}

\begin{remark}[The trivial recurrence]
\label{rem:vacuous}
Every causal sequence model admits a growing-state recurrent form: take
$S_t = x_{1:t}$ and $o_t = M(S_t)$. So the assertion ``a decoder-only
transformer is a recurrent network with growing state'' carries exactly the
content of ``a decoder-only transformer is causal''. Worse for the usual
framing, the trivial state is \emph{smaller} than the KV cache --- $t\log_2 V$
bits against $\Theta(t\,d\,L)$ scalars --- so a transformer does not
\emph{compress} its history into its state; it expands it.
\end{remark}

The content is therefore not existence but cost, and the theorem should be
stated that way.

\begin{theorem}[KV cache as an incrementally maintainable sufficient statistic]
\label{thm:kvform}
Let $M$ be a decoder-only transformer with $L$ layers, width $d$, and $H_{kv}$
key-value heads of dimension $d_h$. Put
$S_t = \big(K^{(1:L)}_{1:t}, V^{(1:L)}_{1:t}\big)$. Then
\begin{enumerate}
  \item[(i)] \emph{(Recurrence)} there exist $f,g$ with $S_t = f(S_{t-1},x_t)$,
        $o_t = g(S_t,x_t)$;
  \item[(ii)] \emph{(Sufficiency)} $o_t$ depends on $x_{1:t-1}$ only through $S_{t-1}$;
  \item[(iii)] \emph{(Incrementality)} $f$ appends $2LH_{kv}d_h$ scalars to
        $S_{t-1}$ and modifies no existing entry; $f$ and $g$ together cost
        $\Theta\!\left(Ld^{2} + L\,t\,H_{kv}d_h\right)$ operations, linear in $t$;
  \item[(iv)] \emph{(Size)} $\dim(S_t) = 2\,t\,L\,H_{kv}d_h$.
\end{enumerate}
The trivial recurrence of Remark~\ref{rem:vacuous} satisfies (i) and (ii) with a
smaller state but costs $\Theta(L t^{2} d)$ per step.
\end{theorem}

\begin{proof}
For (i), causal masking gives that layer-$\ell$ outputs at position $t$ depend on
positions $<t$ only through $K^{(\ell)},V^{(\ell)}$; but the map is not
componentwise, since for $\ell \ge 2$ the pair $(k^{(\ell)}_t, v^{(\ell)}_t)$ is
a function of $\hs{\ell-1}_t$ and hence of $S_{t-1}$. Define $f$ as the $L$-fold
composition that, for $\ell = 1,\dots,L$, forms
$(q,k,v)^{(\ell)}_t$ from $N(\hs{\ell-1}_t)$, appends $(k,v)^{(\ell)}_t$ to
$S^{(\ell)}_{t-1}$, and sets
$\hs{\ell}_t = \hs{\ell-1}_t + \mathrm{Attn}(q^{(\ell)}_t, S^{(\ell)}_t) + \mathrm{MLP}(\cdot)$.
This is a finite composition of maps of $(S_{t-1},x_t)$, giving (i), and (ii) is
immediate from it. For (iii): the projections cost $\Theta(d^2)$ per layer; the
readout at layer $\ell$ is a dot product against $t$ cached keys, costing
$\Theta(t H_{kv} d_h)$; no cached entry is modified because $(k_s,v_s)^{(\ell)}$
depends only on $\hs{\ell-1}_s$ for $s<t$. Summing over $L$ layers gives the
bound. (iv) is counting, and the contrast clause is the cost of re-running the
forward pass on a length-$t$ prefix.
\end{proof}

\begin{remark}
Clause (iii) is the only one distinguishing the KV cache from storing the raw
tokens, and it is what the informal claim omits. Under grouped-query attention
the size in (iv) carries an extra factor $H_{kv}/H$, typically $1/8$ to $1/16$,
which is directly in competition with the capacity question of
Theorem~\ref{thm:recall}. We also note that ``is an RNN'' is an abuse: an RNN
iterates a \emph{single} map on a \emph{fixed} state space, whereas here the
state space grows with $t$ and $f$ is a family of maps with differing domains.
\end{remark}

\subsection{Linear attention}

\begin{theorem}[Linear attention is a bounded-state recurrence]
\label{thm:linattn}
Let $\varphi : \R^{d} \to \R^{d_k}_{\ge 0}$ and set $S_0 = 0 \in \R^{d_k \times d_v}$,
$z_0 = 0 \in \R^{d_k}$. Then attention with kernel
$\exp(\langle q,k\rangle) \mapsto \langle \varphi(q), \varphi(k)\rangle$ satisfies
\begin{equation}
  S_t = S_{t-1} + \varphi(k_t) v_t^{\top},
  \qquad
  z_t = z_{t-1} + \varphi(k_t),
  \qquad
  o_t = \frac{\varphi(q_t)^{\top} S_t}{\varphi(q_t)^{\top} z_t},
  \label{eq:linattn}
\end{equation}
a recurrence in the sense of Definition~\ref{def:recurrent} with bounded state
dimension $d_k d_v + d_k$.
\end{theorem}

\begin{proof}
Unrolling \eqref{eq:linattn} gives $S_t = \sum_{s\le t}\varphi(k_s)v_s^\top$ and
$z_t = \sum_{s \le t}\varphi(k_s)$, so
$o_t = \sum_{s\le t} \langle \varphi(q_t),\varphi(k_s)\rangle v_s \,/\,
       \sum_{s\le t} \langle \varphi(q_t),\varphi(k_s)\rangle$,
which is causal attention with the stated kernel. Non-negativity of $\varphi$
ensures the denominator is a sum of non-negative terms, so it neither vanishes
by cancellation nor changes sign, and $o_t$ is a convex combination of
$\{v_s\}_{s \le t}$.
\end{proof}

\begin{remark}[Gating, and a correction]
\label{rem:gating}
The gated form $S_t = \Diag(a_t)S_{t-1} + \varphi(k_t)v_t^\top$ requires care:
decaying $S$ while leaving $z_t$ undecayed drives $o_t \to 0$, so either $z$ is
gated identically or --- as in practice --- the explicit normaliser is dropped in
favour of output normalisation. Further, the structured-masked-attention duality
of \cite{dao2024mamba2} applies to \emph{scalar} gates $a_t I$: unrolling the
vector-gated recurrence gives attention weight
$\langle \varphi(q_t) \odot A_{t,s}, \varphi(k_s)\rangle$, which factors as a
scalar mask times $\langle\varphi(q_t),\varphi(k_s)\rangle$ if and only if $a_t$
is a multiple of the identity. Vector gating is strictly more general and has no
such dual.
\end{remark}

\subsection{The HMM identity, stated correctly}
\label{sec:state-hmm}

\begin{theorem}[Gated dense recurrence and the forward algorithm]
\label{thm:hmm}
Let $A \in \R^{m\times m}$, $b : \mathcal{O} \to \R^{m}$, and define
\begin{equation}
  s_1 = \pi \odot b(o_1), \qquad
  s_t = \big(A^{\top} s_{t-1}\big) \odot b(o_t) \quad (t \ge 2).
  \label{eq:hmmrec}
\end{equation}
Then:
\begin{enumerate}
  \item[(i)] \emph{(Algebraic identity; no hypotheses)} \eqref{eq:hmmrec} is
        termwise identical to the forward recursion
        $\alpha_t(j) = \big[\sum_i \alpha_{t-1}(i)A_{ij}\big]B_j(o_t)$ with
        $b_j(o) = B_j(o)$.
  \item[(ii)] \emph{(Semantics)} if additionally $\pi$ is a distribution, $A$ is
        row-stochastic and $B_j(\cdot)$ are emission distributions, then
        $s_t(j) = \Prob\!\left[o_{1:t},\, z_t = j\right]$ and
        $\sum_j s_t(j) = \Prob[o_{1:t}]$.
  \item[(iii)] \emph{(Viterbi)} if $A, b \ge 0$, replacing $A^\top s$ by
        $\big(\max_i A_{ij}s(i)\big)_j$ yields
        $s_t(j) = \max_{z_{1:t-1}} \Prob[o_{1:t}, z_{1:t-1}, z_t = j]$; recovering
        the maximising path additionally requires back-pointers.
\end{enumerate}
\end{theorem}

\begin{proof}
(i) Coordinate $j$ of $A^\top s_{t-1}$ is $\sum_i s_{t-1}(i)A_{ij}$, and the
Hadamard product multiplies it by $b_j(o_t)$; this is a rearrangement of a
matrix-vector product and uses no sign or normalisation hypothesis. (ii) By
induction: the base case $s_1(j) = \pi_j B_j(o_1) = \Prob[o_1, z_1 = j]$ uses the
standard convention that $\pi$ is the law of $z_1$; the step is the law of total
probability conditioned on $z_{t-1}$. (iii) The $(\max,\times)$ semiring
distributes over $\R_{\ge 0}$, which is where non-negativity enters and is the
only place it is needed; the induction is otherwise identical.
\end{proof}

\begin{remark}[Which hypotheses do what]
Part (i) needs nothing; row-stochasticity buys only the probabilistic reading in
(ii); non-negativity is load-bearing only for (iii). Note also the base case:
initialising with $s_0 = \pi$ and applying \eqref{eq:hmmrec} at $t=1$ applies one
spurious transition, and is \emph{not} equivalent to \eqref{eq:hmmrec} as stated.
\end{remark}

\subsection{The bridge does not exist}

It is tempting to read Theorem~\ref{thm:hmm} as a constrained special case of
gated linear attention (Theorem~\ref{thm:linattn}), i.e.\ to claim that making a
gated linear RNN's state non-negative and its gate stochastic turns it into an
HMM. It does not, and the obstruction is structural rather than technical.

\begin{theorem}[No embedding of the forward algorithm into gated linear attention]
\label{thm:noembed}
Let $\Phi(S) = \Diag(a)S + C$ be one step of the gated linear-attention
recurrence with any gate $a$ and any input-determined $C$, and let
$\Psi_o(s) = \Diag(b(o))A^{\top}s$ be one step of \eqref{eq:hmmrec}. For a
generic HMM $(A,B)$ with $A$ strictly positive there is no injective linear
$\iota : \R^{m} \to \R^{d_k \times d_v}$ and no choice of gates and inputs
satisfying $\Phi \circ \iota = \iota \circ \Psi_o$ for all $o$.
\end{theorem}

\begin{proof}
Two independent obstructions.

\emph{(a) Support monotonicity.} For any $S,S'$ we have
$\Phi(S) - \Phi(S') = \Diag(a)(S-S')$, since $C$ does not depend on the state.
Hence the set of rows on which two trajectories differ is non-increasing along
the recurrence, for every gate and every input. By contrast, taking
$s - s' = e_r$ gives $\Psi_o(s) - \Psi_o(s') = \big(b_j(o)A_{rj}\big)_j$, whose
support is $\{j : A_{rj} > 0,\, b_j(o) > 0\}$ --- all of $[m]$ for strictly
positive $A$ and $b$. Support strictly grows, so no coordinatewise
identification of the HMM state with rows of $S$ can be equivariant.

\emph{(b) Simultaneous diagonalisability.} Allowing an arbitrary injective
linear $\iota$, the conjugacy $\Phi\circ\iota = \iota\circ\Psi_o$ restricted to
$\operatorname{im}\iota$ states that every $\Psi_o$ is diagonal in the single
basis $\iota(e_1),\dots,\iota(e_m)$. Simultaneously diagonalisable matrices
commute, so this forces $[\Psi_o, \Psi_{o'}] = 0$ for all pairs $o,o'$, which
fails for generic $(A,B)$. Hence no such $\iota$ exists.
\end{proof}

\begin{proposition}[The common generalisation]
\label{prop:affine}
Call $s_t = M(u_t)s_{t-1} + c(u_t)$, $s_t \in \R^{N}$, an \emph{input-controlled
affine recurrence}. Gated linear attention is the case
$M(u_t) = \Diag(a_t)\otimes I_{d_v}$ --- block diagonal, no mixing across the
$d_k$ index --- with $c(u_t) = \mathrm{vec}(\varphi(k_t)v_t^\top)$; the forward
algorithm is the case $M(o_t) = \Diag(b(o_t))A^\top$ dense with $c \equiv 0$.
These are two incomparable cases, not a case and its restriction.
\end{proposition}

\begin{remark}[Consequence for the complaint]
\label{rem:thesis3}
A gate \emph{is} a diagonal matrix; constraining its entries to be non-negative
and to sum to one yields a diagonal stochastic matrix, that is, the identity ---
an HMM whose hidden state never moves. The transition matrix must be
\emph{added} to the architecture, not recovered by constraining something already
present. Diagonality is not incidental: it is precisely what lets
\cite{yang2024gla,dao2024mamba2} evaluate the recurrence as an associative scan.
Architectures whose state transition is dense, such as the delta rule's
$I - \beta_t k_tk_t^\top$ \cite{yang2024delta}, are the ones for which an HMM
reading is even a candidate. The informal thesis --- ``constrain a gated linear
RNN and you get HMM-like decoding'' --- is therefore false as stated, and the
defensible version is that HMM-like decoding requires giving up the diagonal
structure that makes these models fast.
\end{remark}

\subsection{What unbounded state buys}

\begin{definition}[$m$-bit statefulness]
\label{def:mbit}
A sequence model is \emph{$m$-bit stateful} if for every $t$ its state $S_t$
takes values in a set of cardinality at most $2^{m}$, and its emissions after
time $t$ are a deterministic function of $S_t$ and subsequent inputs.
\end{definition}

The finiteness in Definition~\ref{def:mbit} is essential and is exactly what the
informal version of the next result omits: without it the bound is simply false,
since $s_t = s_{t-1}/2 + x_t/2$ stores an unbounded binary expansion in one real
number and can be decoded indefinitely by $x \mapsto \lfloor 2s\rfloor$,
$s \mapsto 2s - \lfloor 2s \rfloor$.

\begin{theorem}[Recall lower bound]
\label{thm:recall}
Let a model be $m$-bit stateful and suppose that for every
$c \in \Vocab^{n}$, when presented with a context containing $c$ at a designated
location followed by a retrieval cue, it emits $c$ exactly. Then
$m \ge n\log_2 V$. If instead it succeeds with probability $1-\delta$ over a
uniform $c$, then $m \ge (1-\delta)\,n\log_2 V - 1$.
\end{theorem}

\begin{proof}
For the exact case, let $S(c)$ be the state reached after the context containing
$c$ and before the cue. Emissions after that point are a function of $S$ and the
cue alone, so $c \mapsto S(c)$ must be injective on $\Vocab^n$; hence
$2^{m} \ge V^{n}$. For the probabilistic case, let $C$ be uniform on $\Vocab^n$
and $\hat C$ the emitted block, a function of $S$. By Fano's inequality
$H(C \mid S) \le 1 + \delta\, n\log_2 V$, while
$H(C\mid S) \ge H(C) - H(S) \ge n\log_2 V - m$; rearranging gives the claim.
\end{proof}

\begin{remark}
The order of quantifiers matters: the block contents must be fixed before the
state forms, and if the retrieval \emph{location} is also adversarial it must be
chosen afterwards, in which case the bound becomes $T\log_2 V$ for a
length-$T$ context. The converse --- that a KV-cache transformer achieves recall
--- is not a corollary of this lower bound but a separate constructive fact
\cite{jelassi2024,arora2023zoology}. Finally, the Fano form predicts
\emph{graceful} degradation past the capacity threshold rather than a sharp
cliff, which is the prediction our window-attention probe should be read against.
\end{remark}
